\documentclass[12pt,a4paper]{article}

% ============================================================
% PACKAGES
% ============================================================
\usepackage[left=1in,right=1in,top=1in,bottom=1in]{geometry}
\usepackage{graphicx}
\usepackage{booktabs}
\usepackage{array}
\usepackage{tabularx}
\usepackage{colortbl}
\usepackage{xcolor}
\usepackage{makecell}
\usepackage{fancyhdr}
\usepackage{titlesec}
\usepackage{amsmath}
\usepackage{tcolorbox}

\definecolor{headerblue}{RGB}{213,220,228}

\newcolumntype{Y}{>{\centering\arraybackslash}X}
\newcolumntype{L}{>{\raggedright\arraybackslash}X}

\setlength{\parindent}{0pt}
\setlength{\parskip}{6pt}

\pagestyle{fancy}
\fancyhf{}
\rhead{SE: Machine Learning Lab}
\lhead{Lab 7}
\rfoot{\thepage}

\begin{document}

% ============================================================
% HEADER
% ============================================================
\renewcommand{\arraystretch}{2.0}
\begin{tabularx}{\textwidth}{|Y|}
\hline
\rowcolor{headerblue} \textbf{\large Department of Computer and Software Engineering} \\
\hline
\textbf{\large SE: Machine Learning} \\
\hline
\end{tabularx}
\renewcommand{\arraystretch}{1.0}

\vspace{12pt}

% ============================================================
% COURSE INFO
% ============================================================
\renewcommand{\arraystretch}{2.0}
\begin{tabularx}{\textwidth}{|L|L|}
\hline
\textbf{Course Instructor:} & \textbf{Date:} \\
\hline
\textbf{Day:} & \textbf{Semester:} \\
\hline
\end{tabularx}
\renewcommand{\arraystretch}{1.0}

\vspace{12pt}

% ============================================================
% TITLE
% ============================================================
\begin{center}
{\Large\bfseries Lab 8: Fraud Detection using Isolation Forest and Gaussian Density}
\end{center}

\vspace{6pt}

% ============================================================
% META TABLE
% ============================================================
\renewcommand{\arraystretch}{1.8}
\begin{tabularx}{\textwidth}{|Y|c|c|c|c|}
\hline
\textbf{Lab Title} & \textbf{CLO} & \textbf{BT} & \textbf{PLO} & \textbf{Weightage} \\
\hline
\makecell{{\small Fraud Detection} \\ {\small using Isolation Forest and Gaussian Density}} & & & & \\
\hline
\end{tabularx}
\renewcommand{\arraystretch}{1.0}

\vspace{6pt}

% ============================================================
% STUDENT TABLE
% ============================================================
\renewcommand{\arraystretch}{2.0}
\begin{tabularx}{\textwidth}{|Y|Y|Y|Y|Y|}
\hline
\textbf{Name} & \textbf{Roll Number} & \textbf{Report /10} & \textbf{Viva /5} & \textbf{Total /15} \\
\hline
 & & & & \\
\hline
\end{tabularx}
\renewcommand{\arraystretch}{1.0}

\vspace{12pt}

\begin{flushright}
Checked on: \_\_\_\_\_\_\_\_\_\_\_\_\_\_\_\_

\vspace{6pt}

Signature: \_\_\_\_\_\_\_\_\_\_\_\_\_\_\_\_
\end{flushright}

\newpage

% ============================================================
% OBJECTIVES
% ============================================================

\section*{Objectives}

\begin{itemize}
\item Understand anomaly detection in real-world fraud scenarios
\item Implement Isolation Forest using sklearn
\item Implement Gaussian density estimation from scratch
\item Visualize anomalous regions
\end{itemize}

% ============================================================
% PROBLEM STATEMENT
% ============================================================

\section*{Problem Statement}

You are designing a fraud detection system for a bank. Each transaction is represented using two features:

\begin{itemize}
\item Transaction Amount
\item Time Gap Since Last Transaction
\end{itemize}

Normal transactions follow consistent patterns, while fraudulent transactions tend to deviate significantly.

Your task is to detect fraudulent transactions using:

\begin{itemize}
\item Isolation Forest (algorithmic approach)
\item Gaussian Density Estimation (probabilistic approach)
\end{itemize}

% ============================================================
% DATASET
% ============================================================

\section*{Dataset Description}

A synthetic dataset is provided:

\begin{itemize}
\item Normal data: Gaussian distribution
\item Fraud data: Random outliers
\item Labels:
    \begin{itemize}
    \item 1 = Normal
    \item -1 = Fraud
    \end{itemize}
\end{itemize}

% ============================================================
% TASKS
% ============================================================

\section*{Part A: Isolation Forest (9 Marks)}

\begin{tcolorbox}
\textbf{Task 1: Train Isolation Forest}

Implement a function to train an Isolation Forest model on the given dataset. 
The model isolates anomalies by recursively splitting the data using random 
feature selection and split values.

\vspace{2cm}

\vspace{2cm}
\end{tcolorbox}

\vspace{6pt}

\begin{tcolorbox}
\textbf{Task 2: Predict Anomalies}

Use the trained model to classify each transaction as normal or anomalous. 
Observe how many points are identified as anomalies and compare with the 
actual distribution.

\vspace{2cm}

\vspace{2cm}
\end{tcolorbox}

\vspace{6pt}

\begin{tcolorbox}
\textbf{Task 3: Analyze Anomaly Scores}

Compute anomaly scores for all data points. These scores are related to 
the path length required to isolate a point in the trees. Points that are 
isolated quickly (short path length) are more likely to be anomalies.

\vspace{2cm}

\vspace{2cm}
\end{tcolorbox}

% ============================================================

\section*{Part B: Gaussian Density Estimation (12 Marks)}

\begin{tcolorbox}
\textbf{Task 4: Estimate Gaussian Parameters}

Compute the mean and covariance matrix of the dataset.

\[
\mu = \frac{1}{N} \sum_{i=1}^{N} x_i
\]

\[
\Sigma = \frac{1}{N} \sum_{i=1}^{N} (x_i - \mu)(x_i - \mu)^T
\]

These parameters represent the center and spread of the data.

\vspace{2cm}

\vspace{2cm}
\end{tcolorbox}

\vspace{6pt}

\begin{tcolorbox}
\textbf{Task 5: Compute Probability Density}

Compute the probability density of each data point using the multivariate 
Gaussian distribution:

\[
p(x) = \frac{1}{(2\pi)^{d/2} |\Sigma|^{1/2}} 
\exp\left(-\frac{1}{2}(x-\mu)^T \Sigma^{-1} (x-\mu)\right)
\]

Points far from the mean will have lower probability values.

\vspace{2cm}

\vspace{2cm}
\end{tcolorbox}

\vspace{6pt}

\begin{tcolorbox}
\textbf{Task 6: Detect Anomalies using Threshold}

Classify transactions based on probability values:

\[
\text{If } p(x) < \epsilon \Rightarrow \text{Anomaly}
\]

\[
\text{Else } \Rightarrow \text{Normal}
\]

Experiment with different values of $\epsilon$ and observe how the classification changes.

\vspace{2cm}

\vspace{2cm}
\end{tcolorbox}

% ============================================================

\section*{Part C: Visualization (6 Marks)}

\begin{tcolorbox}
\textbf{Task 7: Visualize Dataset}

Plot the dataset in a two-dimensional space. Clearly distinguish between 
normal and anomalous transactions using different colors.

\vspace{2cm}

\vspace{2cm}
\end{tcolorbox}

\vspace{6pt}

\begin{tcolorbox}
\textbf{Task 8: Plot Decision Boundaries}

Visualize how both methods (Isolation Forest and Gaussian) separate normal 
data from anomalies. Compare the shape of the regions produced by each method.

\vspace{2cm}

\vspace{2cm}
\end{tcolorbox}

\end{document}